%% Cover letter (draft, for the authors to edit and sign): Paper B to Annals of Global Analysis and
%% Geometry. Build: pdflatex cover-letter-B-AGAG.tex
\documentclass[11pt]{article}
\usepackage[a4paper,top=16mm,bottom=16mm,left=20mm,right=20mm]{geometry}
\usepackage{amsmath,amssymb}
\usepackage[hidelinks]{hyperref}
\setlength{\parindent}{0pt}
\setlength{\parskip}{5pt}
\sloppy
\pagestyle{empty}

\begin{document}

\begin{flushright}
Palaash Gang\\
Indus International School, Pune, India\\
\texttt{palaash.gang@indusschoolpune.com}\\[6pt]
\today
\end{flushright}

The Editors-in-Chief\\
\emph{Annals of Global Analysis and Geometry}

Dear Editors,

On behalf of all six authors (Palaash Gang, Akshaj Agadi, Arjun Veluri, Jerry Wang, Jeremy Barreto and Aarin Chouthaiwale), I submit the manuscript \emph{Finitely many eigenvalues determine the signature of a hyperbolic orbifold} for consideration in the \emph{Annals of Global Analysis and Geometry}. It is posted as arXiv:[PLACEHOLDER: arXiv number of Paper B].

\textbf{What the paper does.} The Laplace spectrum of a closed orientable hyperbolic $2$-orbifold determines its signature, that is, the genus and the orders of the cone points. A measurement or a computation, however, yields only finitely many eigenvalues, each to some accuracy. The paper shows that this suffices, effectively: among orbifolds of area at most $A$, systole at least $\varepsilon$ and cone orders at most $M$, the first $N$ eigenvalues, each known to within $\delta$, determine the signature, with $N$ and $\delta$ given by explicit formulas. This is an effective, signature-level counterpart of the theorem of Buser and Courtois for surfaces. The constants are very large. An a-posteriori test is also given and applied to ten computed spectra; the paper states which of its inputs are estimated rather than proved. Finally, the bound on the cone orders cannot be dropped: the orbifolds $\mathcal O(2,3,m)$ acquire arbitrarily many eigenvalues below any level above $\frac14$.

\textbf{Why AGAG.} The paper belongs to the study of spectral invariants of orbifolds that the journal has published, including Stanhope's spectral bounds on orbifold isotropy (\emph{Ann.\ Global Anal.\ Geom.}\ 27, 2005) and Rossetti, Schueth and Weilandt on isospectral orbifolds with different maximal isotropy orders (34, 2008), both cited in the paper. It also uses the companion paper below, which is submitted to this journal at the same time.

\textbf{Companion manuscripts and use of their results.}
\begin{itemize}\setlength{\itemsep}{2pt}
\item \emph{How much of a hyperbolic orbifold does heat hear?}, by the same authors, submitted to the \emph{Annals of Global Analysis and Geometry} at the same time and cited as arXiv:[PLACEHOLDER: arXiv number of Paper A]. The present paper uses some of its results on the heat expansion and on heat invariants of signatures. These are stated without proof, and Table~1 lists each one with its place in that paper. Nothing else is taken from it. The diameter bound, the effective theorem, the a-posteriori test and the discussion of the hypotheses (Sections~4 to~8) are not contained in it. The two papers share their computed spectra and one code and data deposit.
\item \emph{Triples with equal sum and equal reciprocal sum}, by the same authors, submitted to \emph{Integers} (arXiv:[PLACEHOLDER: arXiv number of the arithmetic note]). The present paper does not use it.
\end{itemize}
Because the present paper relies on the first companion, we respectfully suggest that the same editor handle both manuscripts.

\textbf{Code and data.} The code and data are publicly archived at Zenodo, with all six authors as creators: [PLACEHOLDER: Zenodo DOI, to be minted at submission]. One command reruns every check, and every constant of the tables of explicit bounds is recomputed at $40$ digits.

The manuscript is not under consideration elsewhere, and all authors have approved its submission. We would be grateful if you would consider it for publication.

Yours sincerely,

\vspace{10mm}
Palaash Gang (corresponding author), on behalf of all authors

\end{document}
